% Part: first-order-logic
% Chapter: tableaux
% Section: proof-theoretic-notions

\documentclass[../../../include/open-logic-section]{subfiles}

\begin{document}

\iftag{FOL}
      {\olfileid[fr]{fol}{tab}{ptn}}
      {\olfileid[fr]{pl}{tab}{ptn}}
      
\olsection{Notions de théorie de la démonstration}

\begin{editorial}
Cette section rassemble les définitions de la relation de
!!{derivability} et de la cohérence pour les tableaux.
\end{editorial}

\begin{explain}
Nous avons défini plusieurs notions sémantiques importantes : validité,
conséquence logique et satisfaisabilité. Définissons maintenant les
\emph{notions de théorie de la démonstration} correspondantes. Elles ne
font pas appel à la satisfaction d'!!{sentence}s dans des
\iftag{FOL}{!!{structure}s}{!!{valuation}s}, mais à l'existence de
certains !!{tableau}x fermés. La coïncidence de ces
notions constitue une découverte importante : c'est le contenu des
\emph{théorèmes de correction et de complétude}.
\end{explain}


\begin{defn}[Théorèmes]
Un !!{sentence}~$!A$ est un \emph{théorème} s'il existe un
!!{tableau} fermé pour~$\sFmla{\False}{!A}$. On écrit $\Proves !A$
si $!A$ est un théorème, et $\Proves/ !A$ dans le cas contraire.
\end{defn}

\begin{defn}[!!^{derivability}]
!!^a{sentence}~$!A$ est \emph{!!{derivable} à partir} d'un ensemble
d'!!{sentence}s~$\Gamma$, ce que l'on note $\Gamma \Proves !A$,
si et seulement s'il existe une partie finie
$\{!B_1, \dots, !B_n\} \subseteq \Gamma$ et un !!{tableau} fermé
pour l'ensemble
\[
\{
  \sFmla{\False}{!A},
  \sFmla{\True}{!B_1}, \dots,
  \sFmla{\True}{!B_n}
\}.
\]
Si $!A$ n'est pas !!{derivable} à partir de $\Gamma$, on écrit
$\Gamma \Proves/ !A$.
\end{defn}

\begin{defn}[Cohérence]
Un ensemble d'!!{sentence}s~$\Gamma$ est \emph{incohérent} si et
seulement s'il existe une partie finie
$\{!B_1, \dots, !B_n\} \subseteq \Gamma$ et un !!{tableau} fermé
pour l'ensemble
\[
\{
  \sFmla{\True}{!B_1}, \dots,
  \sFmla{\True}{!B_n}  
\}.
\]
Si $\Gamma$ n'est pas incohérent, on dit qu'il est \emph{cohérent}.
\end{defn}

\begin{prop}[Réflexivité]
\ollabel{prop:reflexivity}
Si $!A \in \Gamma$, alors $\Gamma \Proves !A$.
\end{prop}

\begin{proof}
Si $!A \in \Gamma$, alors $\{!A\}$ est une partie finie de~$\Gamma$
et le !!{tableau}
\begin{oltableau}
  [\sFmla{\False}{\formula{A}}, just = \TAss
    [\sFmla{\True}{\formula{A}}, just = \TAss,close]
  ]
\end{oltableau}
est fermé.
\end{proof}
  

\begin{prop}[Monotonie]
\ollabel{prop:monotonicity}
Si $\Gamma \subseteq \Delta$ et $\Gamma \Proves !A$, alors
$\Delta \Proves !A$.
\end{prop}

\begin{proof}
Toute partie finie de~$\Gamma$ est aussi une partie finie de~$\Delta$.
\end{proof}

\begin{prop}[Transitivité]
\ollabel{prop:transitivity}
Si $\Gamma \Proves !A$ et $\{!A\} \cup
\Delta \Proves !B$, alors $\Gamma \cup \Delta \Proves !B$.
\end{prop}

\begin{proof}
Si $\{!A\} \cup \Delta \Proves !B$, il existe une partie finie
$\Delta_0 = \{!C_1, \dots, !C_n\} \subseteq \Delta$ telle que
\begin{align*}
\{\sFmla{\False}{!B}, & \sFmla{\True}{!A}, \sFmla{\True}{!C_1},
\dots, \sFmla{\True}{!C_n}\}
\intertext{possède un !!{tableau} fermé. Si $\Gamma \Proves !A$,
  il existe une partie finie
  $\{!D_1, \dots, !D_m\} \subseteq \Gamma$ telle que}
\{\sFmla{\False}{!A}, & \sFmla{\True}{!D_1},
\dots, \sFmla{\True}{!D_m}\}
\end{align*}
possède un !!{tableau} fermé.

Considérons maintenant le !!{tableau} ayant pour hypothèses
\[
\sFmla{\False}{!B},
\sFmla{\True}{!C_1}, \dots, \sFmla{\True}{!C_n},
\sFmla{\True}{!D_1}, \dots, \sFmla{\True}{!D_m}.
\]
Appliquons la règle \Cut{} à~$!A$. Elle produit deux branches : l'une
contient $\sFmla{\True}{!A}$, l'autre $\sFmla{\False}{!A}$. Sur la
première branche, toutes les formules de
\[
\{\sFmla{\False}{!B}, \sFmla{\True}{!A},
\sFmla{\True}{!C_1}, \dots, \sFmla{\True}{!C_n}\}
\]
sont disponibles. Puisqu'il existe un !!{tableau} fermé pour ces
hypothèses, nous pouvons le greffer à cette branche; toute branche qui
passe par $\sFmla{\True}{!A}$ se ferme. Sur l'autre branche, toutes les
formules de
\[
\{\sFmla{\False}{!A}, \sFmla{\True}{!D_1}, \dots,
\sFmla{\True}{!D_m}\}
\]
sont disponibles; nous pouvons donc compléter cette branche elle aussi
et obtenir un !!{tableau} fermé. Cela montre que
$\Gamma \cup \Delta \Proves !B$.
\end{proof}

\begin{editorial}
Dans le paragraphe introductif, la mention des structures dans la
source est adaptée en valuations pour la présentation propositionnelle.
Dans la preuve de transitivité, la source écrit par erreur
$!D_1, \dots, !D_m \subseteq \Gamma$. Les accolades rétablissent la
partie finie $\{!D_1, \dots, !D_m\} \subseteq \Gamma$ exigée par la
définition de la dérivabilité.
\end{editorial}

Remarquons que cela signifie en particulier que, si
$\Gamma \Proves !A$ et $!A \Proves !B$, alors $\Gamma \Proves !B$.
Il en résulte aussi que, si $!A_1, \dots, !A_n \Proves !B$ et
$\Gamma \Proves !A_i$ pour chaque~$i$, alors $\Gamma \Proves !B$.

\begin{prop}
\ollabel{prop:incons}
$\Gamma$ est incohérent si et seulement si $\Gamma \Proves !A$
pour tout !!{sentence}~$!A$.
\end{prop}

\begin{proof}
Exercice.
\end{proof}

\tagprob{FOL}
\begin{prob}
Démontrer la \olref[fol][tab][ptn]{prop:incons}.
\end{prob}
\tagendprob

\tagprob{notFOL}
\begin{prob}
Démontrer la \olref[pl][tab][ptn]{prop:incons}.
\end{prob}
\tagendprob

\begin{prop}[Compacité]
  \ollabel{prop:proves-compact}
  \begin{enumerate}
  \item Si $\Gamma \Proves !A$, il existe une partie finie $\Gamma_0
    \subseteq \Gamma$ telle que $\Gamma_0 \Proves !A$.
  \item Si toute partie finie de~$\Gamma$ est cohérente,
    alors $\Gamma$ est cohérent.
  \end{enumerate}
\end{prop}

\begin{proof}
  \begin{enumerate}
    \item Si $\Gamma \Proves !A$, il existe une partie finie
      de~$\Gamma$, notée
      $\Gamma_0 = \{!B_1, \dots, !B_n\}$, et un !!{tableau} fermé pour
      \[
      \{\sFmla{\False}{!A}, \sFmla{\True}{!B_1}, \dots, \sFmla{\True}{!B_n}\}
      \]
      Ce !!{tableau} montre aussi que $\Gamma_0 \Proves !A$.
    \item Si $\Gamma$ est incohérent, il existe une partie finie
      de~$\Gamma$, notée
      $\Gamma_0 = \{!B_1, \dots, !B_n\}$, et un !!{tableau} fermé pour
      \[
      \{\sFmla{\True}{!B_1}, \dots, \sFmla{\True}{!B_n}\}
      \]
      Ce !!{tableau} fermé montre que $\Gamma_0$ est incohérent.
  \end{enumerate}
\end{proof}

\end{document}
